% from main.tex
%%%%%%%%%%  scan idx 275  |  folio 262  %%%%%%%%%%
condition a) (resp. a'), etc.) relative à \(E_x\) et \(\Mod(\mathcal O_{X,x})\). Alors, si \(X\) a assez de points, la condition a) (resp. a'), b)) équivaut à la condition «~pour tout point \(x\) de \(X\), \(\mathrm a_x)\) (resp. \(\mathrm a'_x)\), \(\mathrm b_x))\)~», et équivaut à la condition «~pour tout point \(x\) de \(X\), \(\mathrm b'_x)\)~» si \(E\) est pseudo-cohérent ; en outre la condition a') équivaut à la condition analogue où l'on remplace «~pour tout \(M\in\Ob\Mod(X)\)~» par «~pour tout \(\mathcal O_X\)-Module conoyau d'un homomorphisme \(\mathcal O_{U,X}^p\longrightarrow\mathcal O_{U,X}^q\) où \(U\) est un objet de \(X\)~».

(iii) Le critère (EGA III 7.4.2) s'étend mot pour mot à notre situation : si \(E\) est à composantes plates, la condition a) signifie simplement que \(\Coker(E^p\longrightarrow E^{p+1})\) est plat, donc localement facteur direct d'un Module libre de type fini si \(E^p\) et \(E^{p+1}\) sont supposés être en outre de présentation finie (I 5.8.3). On en déduit en particulier que, si \(E\) est parfait, alors \(E\) vérifie la condition a) relative à l'entier \(p\) si et seulement si \(E^{\vee}\) vérifie la condition a) relative à l'entier \(-p-1\).

(iv) Soient \(X\) un schéma quasi-séparé, \(\mathcal O_X\) le faisceau structural, \(E\in\Ob D^-(X)\) un complexe à cohomologie localement bornée et quasi-cohérente. Notons \(\mathrm a_{\qcoh})\) (resp. \(\mathrm a'_{\qcoh})\), etc.) la condition analogue à a) (resp. a'), etc). où l'on remplace \(\Mod(X)\) par la catégorie \(\Qcoh(X)\) des \(\mathcal O_X\)-Modules quasi-cohérents. Alors on a
\[
 \mathrm a)\Longleftrightarrow\mathrm a_{\qcoh})\Longleftrightarrow\mathrm a'_{\qcoh})\Longleftrightarrow\mathrm b_{\qcoh})\Longleftarrow\mathrm b'_{\qcoh}),
\]
et les conditions précédentes sont équivalentes si \(E\) est pseudo-cohérent.

En effet, comme \(X\) est quasi-séparé, tout faisceau quasi-cohérent sur un ouvert affine de \(X\) est induit par un faisceau quasi-cohérent sur \(X\) (EGA I 9.4.2 et IV 1.2.7). Il en résulte que la condition \(\mathrm a_{\qcoh})\) (resp. \(\mathrm a'_{\qcoh})\), \(\mathrm b_{\qcoh})\), \(\mathrm b'_{\qcoh}))\) est vérifiée si et seulement si elle l'est au-dessus de tout ouvert affine. Pour prouver les implications ci-dessus, on peut donc supposer \(X\) affine. Alors (II 3.5) \(E\) est isomorphe dans \(D(X)\) à un complexe de Modules quasi-cohérents, et les implications
\[
 \mathrm a_{\qcoh})\Longleftrightarrow\mathrm a'_{\qcoh})\Longleftrightarrow\mathrm b_{\qcoh})\Longleftarrow\mathrm b'_{\qcoh})
\]
